\section{Detailed numerical verification}
\label{sec:supp-numerics}

The local laws of Section~\ref{sec:rates} predict constants, not merely exponents, so
they can be falsified numerically.  We report closed-form instances in which
the displayed constants are available in exact arithmetic.  The
\(2\times2\) checks evaluate the smoothed scalar root of
\eqref{eq:smooth-root-main} in high-precision decimal arithmetic through the
closed principal-square-root formula for
\(Z_\delta(A)=A(A^\top A+\delta^2I)^{-1/2}\); the diagonal corank-two
instance is evaluated entrywise.  The roots are located by
bisection on the strictly increasing function \(h_\delta\); assertion-based
scripts accompany the paper.

For the first group, \(\Theta=e_1e_1^\top\) and
\(G=\bigl(\begin{smallmatrix}0&b_1\\b_2&c\end{smallmatrix}\bigr)\).

\begin{center}
\begin{tabular}{@{}llll@{}}
\toprule
Regime & Instance and predicted constant & Smoothing & Computed \\
\midrule
Theorem~\ref{thm:rate-regular} &
\(b_1=b_2=1,\ c=\tfrac12\);\ \(\lambda^\star=\tfrac12\), \(\kappa=\tfrac12\) &
\(\delta=10^{-4}\) & \(-1.777777665\)\\
\(O(\delta^2)\) &
\((\lambda_\delta-\lambda^\star)/\delta^2\to-b^\star/\kappa=-\tfrac{16}{9}\)
&\(\delta=10^{-5}\) & \(-1.777777777\)\\
\addlinespace
Corollary~\ref{thm:rate-strict} &
\(b_1=b_2=\tfrac12,\ c=1\);\ \(\lambda^\star=\tfrac14\),
\(a=\tfrac15\), \(\beta=\tfrac45\) &
\(\delta=10^{-6}\) & \(-0.3227484177\)\\
\(\Theta(\delta)\) &
\((\lambda_\delta-\lambda^\star)/\delta\to t_0=-\tfrac{5}{4\sqrt{15}}\)
& \(\delta=10^{-8}\) & \(-0.3227486102\)\\
\addlinespace
Theorem~\ref{thm:rate-boundary} &
\(b_1=b_2=c=1\);\ \(\lambda^\star=1\), \(a=\beta=\tfrac12\), \(\kappa=\tfrac12\) &
\(\delta=10^{-9}\) & \(1.259920256\)\\
\(\Theta(\delta^{2/3})\) &
\((\lambda^\star-\lambda_\delta)/\delta^{2/3}\to(2\beta\kappa)^{-1/3}=2^{1/3}\)
& \(\delta=10^{-12}\) & \(1.259921042\)\\
\bottomrule
\end{tabular}

\smallskip
{\small
Exact values: \(-16/9=-1.777777778\),
\(-5/(4\sqrt{15})=-0.3227486122\),
\(2^{1/3}=1.259921050\).}
\end{center}

For the same boundary data, standard Huber gives
\((\lambda^\star-\lambda_\delta^{\rm hub})/\delta=1.9999999656\) at
\(\delta=10^{-8}\), versus \(1/\beta=2\).  The scaled-Huber regression in
the code also confirms that moving the first saturation threshold to
\(L=1/2\) changes the constant to \(L/\beta=1\), as required by
\eqref{eq:saturated-boundary-rate-main}.

The preceding instances are \(2\times2\).  To exercise the general
completion and arbitrary-corank strict law, take
\[
 A^\star=\diag(1,1,0,0),\qquad
 \Theta=\diag\!\left(\tfrac38,\tfrac38,1,\tfrac12\right),\qquad
 \lambda^\star=1,\qquad G=A^\star-\Theta.
\]
Then \(a=3/4\), \(C=\diag(1,1/2)\), and \(\beta=3/2\).  Hard
water-filling gives
\(\eta_F=3/5\),
\(K_F=-\diag(3/5,3/10)\), and
\(\|K_F\|_F=\sqrt{0.45}\).  The smooth-center equation gives
\[
 \eta_R=0.715080068980,\qquad
 K_R=-\diag(0.581665969555,0.336668060890),
\]
with \(\|K_F-K_R\|_F=0.040996138377\).  Writing
\(\Phi_F^\star\) and \(\Phi_R^\star\) for the two corresponding completed
certificates, the actual smoothing path gives
\begin{center}
\begin{tabular}{@{}llll@{}}
\toprule
\(\delta\) & \(\lambda_\delta-\lambda^\star\) &
\(\|\Phi_\delta-\Phi_R^\star\|_F\) &
\(\|\Phi_\delta-\Phi_F^\star\|_F\)\\
\midrule
\(10^{-2}\) & \(-7.150296\!\times\!10^{-3}\) & \(7.896\!\times\!10^{-5}\) & \(4.0990\!\times\!10^{-2}\)\\
\(10^{-3}\) & \(-7.150796\!\times\!10^{-4}\) & \(7.859\!\times\!10^{-7}\) & \(4.0996\!\times\!10^{-2}\)\\
\(10^{-4}\) & \(-7.150801\!\times\!10^{-5}\) & \(7.855\!\times\!10^{-9}\) & \(4.0996\!\times\!10^{-2}\)\\
\bottomrule
\end{tabular}
\end{center}
Thus the path converges to the smooth rather than the minimum-Frobenius
center, while
\((\lambda_\delta-\lambda^\star)/\delta\to-\eta_R\), as predicted by
Theorems~\ref{thm:general-fenchel-center} and~\ref{thm:rate-strict-general}
and Corollary~\ref{cor:center-agreement-main}.
On the same line, the standard Huber path gives slope
\(-\eta_F=-3/5\) and converges to \(\Phi_F^\star\) to below
\(10^{-11}\) at \(\delta=10^{-5}\), verifying
Proposition~\ref{prop:smoothing-family-main} on the same nonsingleton slice.

The preceding corank-two line is diagonal: the relevant singular dyads are
fixed.  In the notation of Proposition~\ref{prop:strict-constants-closed-main},
both \(P_{\rm big}(s)\) and \(\widetilde B(s)\) are constant.  Hence
\(\partial_s\mathcal W=0\), which proves \(t_1=0\) and
\(\mathcal M_R=0\) rather than merely observing their cancellation.  To test a nonzero first-order
coefficient without random data, take
\[
 A^\star=\diag(1,3/5,0,0),\qquad \lambda^\star=1,
\]
\[
 \Theta=
 \begin{pmatrix}
  1/6&0&2/5&0\\
  0&1/6&0&3/10\\
  1/5&0&2/3&0\\
  0&1/10&0&1/3
 \end{pmatrix},
 \qquad G=A^\star-\Theta.
\]
In the displayed coordinates, \(a=1/3\),
\(C=\diag(2/3,1/3)\), \(\beta=1\), and
\(\eta_R=0.644812526657\).  At \(\delta=10^{-5}\), the extended-precision
multiplier ratio is \(-0.644805516183\), while the deterministic direction
ratio is \(0.671768254\),
respectively, already exhibiting a nonzero linear direction term.  Direct evaluation of
\eqref{eq:t1-closed-main}--\eqref{eq:MR-closed-main} gives
\[
 t_1=0.7010530600,
 \qquad \|\mathcal M_R\|_F=0.6717753807.
\]
At \(\delta=10^{-6}\), the observed direction ratio is
\(0.6717760680\), and the cosine between
\((\Phi_\delta-\Phi_R^\star)/\delta\) and \(\mathcal M_R\) equals one to
ten displayed digits.  The corresponding 50-digit finite-scale estimate of
\(t_1\) at \(\delta=10^{-5}\) is \(0.7010473954\).  Thus the table verifies the
predicted vector coefficient, not only its exponent.  This demonstrates
sharpness for one fixed instance; no generic measure statement is inferred.

The compact-root branch of
Corollary~\ref{cor:compact-root-superconvergence-main} is sharp as well.
For \(A^\star=\diag(2,1,0)\), \(\Theta=\diag(1,-1,1)\), and
\(G=A^\star-\Theta\), one has \(a=0\), \(C=(1)\),
\(J_\star=\diag(1/4,1,0)\), and \(\chi_\star=-3/4\).  Thus the predicted
coefficients are \(-3/8\) and
\(\|\mathcal N_0\|_F=\sqrt{26}/8=0.6373774392\ldots\).
At \(\delta=10^{-4}\), high-precision bisection gives the normalized pair
\((-0.3749999965,\,0.6373774341)\); the assertion script checks monotone
convergence over four scales.  Thus, even when the compact-polar equation
already has a root, the nearby smooth path has nonzero errors beginning at
cubic and quadratic order.

The doubly cancelled branch in
\eqref{eq:compact-root-quintic-main} uses
\[
 A^\star=\diag(3,2,1,0),\qquad \lambda^\star=1,
\]
\[
 \Theta=\begin{pmatrix}
  27/10&0&0&3/10\\
  0&-16/5&0&1/10\\
  0&0&1/2&1/4\\
  1/5&-3/20&1/10&1
 \end{pmatrix},\qquad G=A^\star-\Theta.
\]
Here \(a=0\), \(C=(1)\), \(d_0=1\),
\(\chi_\star=0\), and \(\chi_2=1/3\).  Hence the predicted quintic
coefficient is \(-1/8\), and the quartic direction coefficient after
subtracting \(-\delta^2J_\star/2\) is
\[
 \diag\!\left(\frac1{216},\frac3{128},\frac38,-\frac18\right).
\]
Extended-precision bisection gives
\begin{center}
\begin{tabular}{@{}lllll@{}}
\toprule
\(\delta\) & \(10^{-2}\) & \(10^{-3}\) & \(10^{-4}\) & \(10^{-5}\)\\
\midrule
\((\lambda_\delta-\lambda^\star)/\delta^5\) &
\(-0.1248207361\)&\(-0.1249833498\)&\(-0.1249983476\)&\(-0.1249998349\)\\
\bottomrule
\end{tabular}
\end{center}
At \(\delta=10^{-4}\), the diagonal of the normalized quartic direction is
\[
 (0.004629629625,,0.023437499951,,0.374999996875,
   -0.124998347562),
\]
while the Frobenius norm of its off-diagonal part is below
\(4\times10^{-6}\).  The high-precision assertion script checks all four
scales and the full matrix coefficient.

For Theorem~\ref{thm:crossover-main} we use the one-parameter family
\(b_1=b_2=1\), \(c=1+\varepsilon\), for which
\(\beta_0=\kappa_0=\tfrac12\) and
\(\Delta_\varepsilon=(c^2-1)/(1+c^2)\).  Choosing \(\varepsilon=
\varepsilon(\delta)\) so that \(\Delta_\delta=\zeta\delta^{2/3}\) places the
family in the critical window \eqref{eq:crossover-critical-main}, whose
limit solves \(\tfrac12y^3+\zeta y^2=1\).

\begin{center}
\begin{tabular}{@{}lllll@{}}
\toprule
\(\zeta\) & \(0\) & \(1\) & \(3\) & \(10\)\\
\midrule
predicted \(y_\zeta\) &
\(1.259921050\)&\(0.839286755\)&\(0.552474611\)&\(0.313775961\)\\
computed \(x_\delta/\delta^{2/3}\), \(\delta=10^{-9}\) &
\(1.259920256\)&\(0.839284833\)&\(0.552471210\)&\(0.313769661\)\\
\bottomrule
\end{tabular}
\end{center}

For a valid strict-side moving path, take
\(\Delta_\delta=\delta^{1/3}\); then
\(\Delta_\delta\to0\) while
\(\Delta_\delta/\delta^{2/3}\to\infty\).  The normalization in
\eqref{eq:crossover-inner-main}, whose limit is one for this family, gives
\begin{center}
\begin{tabular}{@{}lllll@{}}
\toprule
\(\delta\) & \(10^{-6}\) & \(10^{-9}\) & \(10^{-12}\) & \(10^{-15}\)\\
\midrule
\(x_\delta\sqrt{\Delta_\delta}/\delta\) &
\(0.979950782\)&\(0.997994100\)&\(0.999799770\)&\(0.999979992\)\\
\bottomrule
\end{tabular}
\end{center}
This path respects the moving-family quantifier.  By contrast, fixing
\(\varepsilon\) and sending \(\delta\) to zero tests the fixed strict-kink
law, not the crossover theorem.

Finally, the oracle returns a matrix direction, not only a multiplier.  At
the boundary instance \(b_1=b_2=c=1\), direct differentiation gives
\[
 \|\mathcal M_0\|_F=\sqrt{3/2},
 \qquad
 \frac{\|\Phi_\delta-\Phi^\star\|_F}{x_\delta}
 \longrightarrow\sqrt{3/2}.
\]
At \(\delta=10^{-15}\), the computed ratio is
\(1.224744871327\), versus
\(\sqrt{3/2}=1.224744871392\).  The same verification includes critical,
strict-side, boundary, and deliberately oscillatory paths for the unified
ratio \(x_\delta/\widehat x_\delta\).

We emphasize what these tables do and do not support.  They confirm the
exponents and leading constants of the local laws on instances where those
constants are known in closed form.  They are not evidence about how often
rank-deficient crossings occur in any application, and no such claim is made.
